\section{Configuration Management: Source Listings}\label{sec:appendix-configuration-management-listings}

\subsection{NSAK Core Configuration}\label{subsec:appendix-nsak-core-configuration}

\begin{listing}[H]
    \inputminted{python}{../../../src/nsak/core/settings.py}
    \caption{NSAK Static Configuration}
    \label{lst:nsak-settings}
\end{listing}

\begin{listing}[H]
    \inputminted{python}{../../../src/nsak/core/configuration/configuration.py}
    \caption{NSAK Runtime Configuration}
    \label{lst:nsak-config}
\end{listing}

\begin{listing}[H]
    \inputminted{python}{../../../src/nsak/cli/__init__.py}
    \caption{NSAK CLI Initialisation}
    \label{lst:nsak-cli-init}
\end{listing}

\begin{listing}[H]
    \begin{minted}{yaml}
debug: true
device:
  id: bananapi_r4
  name: BananaPI R4
  path: <base-path>/lib/devices/bananapi_r4
  author: vonal3@bfh.ch
  repository: <repository-url>
  configuration:
    network:
      ethernets:
        lan1@eth0:
          addresses:
            10.10.10.30/24:
              is_target: true
              is_management: false
    \end{minted}
    \caption{Persisted Runtime Configuration (run/config.yaml)}
    \label{lst:nsak-config-yaml}
\end{listing}

\begin{listing}[H]
    \begin{minted}{bash}
# Show usage, options and subcommands for the device resource
nsak device --list

# List all devices found in the resource library
nsak device list

# Output:
❯ nsak device list
ID           NAME         PATH
bananapi_r4  BananaPI R4  <repository>/lib/devices/bananapi_r4
    \end{minted}
    \caption{NSAK CLI: List Devices}
    \label{lst:nsak-cli-list-devices}
\end{listing}

\subsection{Device Resource Management}\label{subsec:appendix-device-resource-management}

\begin{listing}[H]
    \begin{minted}{yaml}
devices:
    <str:resource-id>:
        id: <str:device-id>
        name: <str:device-name>
        description: <str:device-description>
        author: <str:author-email>
        repository: <str:resource-repository>
    \end{minted}
    \caption{Device Resource YAML Specification}
    \label{lst:device-resource-yaml-specification}
\end{listing}

\begin{listing}[H]
    \begin{minted}{yaml}
devices:
    <str:resource-id>:
        # ...
        configuration:
          network:
            ethernets:
              <str:intreface-name>:
                addresses:
                  <str:cidr>:
                    is_target: <boolean>
                    is_management: <boolean>
    \end{minted}
    \caption{Device Resource Configuration YAML Specification}
    \label{lst:device-resource-configuration-yaml-specification}
\end{listing}

\begin{listing}[H]
    \inputminted[firstline=20, lastline=87]{python}{../../../src/nsak/core/device/device_loader.py}
    \caption{DeviceLoader Configuration Parsing}
    \label{lst:device-loader-configuration-parsing}
\end{listing}

\begin{listing}[H]
    \inputminted[firstline=12, lastline=30]{python}{../../../src/nsak/core/network/configuration.py}
    \inputminted[firstline=14, lastline=39]{python}{../../../src/nsak/core/device/device.py}
    \caption{Device Configuration Datastructures}
    \label{lst:device-configuration-datastructures}
\end{listing}

\begin{listing}[H]
    \begin{minted}{bash}
# List all subcommands for the Device resource
nsak device --help

# List all available devices
nsak device list

# Show the device details and configuration
nsak device show <str:device-id>

# Load a device into in the NSAK configuration
nsak device load <str:device-id>

# Show the currently loaded device stored in the NSAK configuration
nsak device loaded

# Reset the loaded device
nsak device unload
    \end{minted}
    \caption{Device Configuration CLI Management}
    \label{lst:device-configuration-cli-management}
\end{listing}

\subsection{Resource YAML Specifications}\label{subsec:appendix-resource-yaml-specifications}

\begin{listing}[H]
    \begin{minted}{yaml}
# Resource type version was removed, as it was not used anyway

# scenarios|drills|environments|devices
<str:resource-type-key>:
    # unique resource identifier
    <str:resource-id>:
        # The "metadata" key is removed and the children are now nested directly under the resource
        name: <str:name>
        description: <str:description>
        author: <str:email>
        repository: <str:link>
        # ... resource type specific properties, which previously were at the top level
    \end{minted}
    \caption{Mergeable Resource YAML Specification}
    \label{lst:mergeable-resource-yaml-specification}
\end{listing}

\begin{listing}[H]
    \begin{minted}{yaml}
<str:resource-type-key>:
    <str:resource-id>:
        # Unchanged properties...

        interface:
            arguments:
                <str:argument-name>:
                    type: <str:type-annotation>
                    default: <any:default-value>
            return_type: <str:type-annotation>
    \end{minted}
    \caption{Updated Interface Specification}
    \label{lst:updated-interface-specification}
\end{listing}

\begin{listing}[H]
    \inputminted{yaml}{../../../lib/scenarios/mitm/scenario.yaml}
    \caption{mitm Scenario YAML}
    \label{lst:mitm-scenario-yaml}
\end{listing}
